\documentclass[12pt,a4paper]{article}
\usepackage[utf8]{inputenc}
\usepackage{amsmath,amssymb,amsthm}
\usepackage[margin=1in]{geometry}

\newtheorem{theorem}{Theorem}
\newtheorem{lemma}[theorem]{Lemma}
\newtheorem{proposition}[theorem]{Proposition}
\newtheorem{corollary}[theorem]{Corollary}
\theoremstyle{definition}
\newtheorem{definition}[theorem]{Definition}

\title{Problem 16: Power Digit Sum}
\author{Project Euler Solutions}
\date{}

\begin{document}
\maketitle

\section{Problem Statement}
Let $S(m)$ denote the sum of the decimal digits of a positive integer~$m$. Compute~$S(2^{1000})$.

\section{Formal Development}

\begin{definition}
For a positive integer $m$ with base-10 representation $m = \sum_{k=0}^{d-1} a_k \cdot 10^k$,
where $0 \le a_k \le 9$ and $a_{d-1} \ne 0$, the \emph{digit sum} is $S(m) = \sum_{k=0}^{d-1} a_k$,
and the \emph{digit length} is $d(m) = \lfloor \log_{10} m \rfloor + 1$.
\end{definition}

\begin{theorem}[Digit length of $2^n$]
For all $n \ge 0$,
\[
d(2^n) = \lfloor n \log_{10} 2 \rfloor + 1.
\]
\end{theorem}

\begin{proof}
$d(2^n) = \lfloor \log_{10}(2^n) \rfloor + 1 = \lfloor n \log_{10} 2 \rfloor + 1.$
\end{proof}

\begin{corollary}
$d(2^{1000}) = \lfloor 1000 \cdot 0.30103\ldots \rfloor + 1 = 302$.
\end{corollary}

\begin{theorem}[Digit sum congruence]
For every positive integer $m$, $S(m) \equiv m \pmod{9}$.
\end{theorem}

\begin{proof}
Write $m = \sum_{k=0}^{d-1} a_k \cdot 10^k$. Since $10 \equiv 1 \pmod{9}$, we have
$10^k \equiv 1 \pmod{9}$ for all $k \ge 0$, so $m \equiv \sum a_k = S(m) \pmod{9}$.
\end{proof}

\begin{lemma}[Order of 2 modulo 9]
$\mathrm{ord}_9(2) = 6$.
\end{lemma}

\begin{proof}
Direct computation: $2^1 \equiv 2$, $2^2 \equiv 4$, $2^3 \equiv 8$, $2^4 \equiv 7$,
$2^5 \equiv 5$, $2^6 \equiv 1 \pmod{9}$, and no smaller positive exponent yields~$1$.
\end{proof}

\begin{proposition}
$2^{1000} \equiv 7 \pmod{9}$.
\end{proposition}

\begin{proof}
Since $1000 = 166 \cdot 6 + 4$ and $\mathrm{ord}_9(2) = 6$, we have
$2^{1000} \equiv 2^4 = 16 \equiv 7 \pmod{9}$.
\end{proof}

\begin{theorem}[Bounds]
$1 \le S(2^n) \le 9\bigl(\lfloor n \log_{10} 2 \rfloor + 1\bigr)$ for all $n \ge 1$.
\end{theorem}

\begin{proof}
The leading digit of $2^n$ is at least~$1$, giving the lower bound. Each of the
$d(2^n)$ digits is at most~$9$, giving the upper bound.
\end{proof}

\paragraph{Verification.}
The answer $S(2^{1000}) = 1366$ satisfies $1366 = 151 \cdot 9 + 7 \equiv 7 \pmod{9}$,
consistent with $2^{1000} \equiv 7 \pmod{9}$.

\section{Complexity}

\begin{proposition}
The iterative doubling algorithm runs in $\Theta(n^2)$ time and $\Theta(n)$ space.
\end{proposition}

\begin{proof}
At iteration~$i$, the digit array has length $\Theta(i)$. The total work is
$\sum_{i=1}^{n} \Theta(i) = \Theta(n^2)$. The array stores $O(n)$ digits.
\end{proof}

\section{Answer}

\[
\boxed{1366}
\]

\end{document}
